\begin{table}[tp]
\centering
\caption{Heavy-tailed composite acceleration under \href{https://github.com/k3nfalt/longtails/blob/main/LongTails/Setups/StochasticCompositeGeneralizedSmooth.lean}{StochasticCompositeGeneralizedSmooth}.}\label{tab:StochasticCompositeGeneralizedSmooth}
\parbox{\linewidth}{\footnotesize On a second-countable real Hilbert space, minimize $F=f+\psi$ over a measurable convex set $Q$. The convex $C^1$ function $f$ satisfies $\|\nabla f(x)-\nabla f(y)\|\le(L_0+L_1\sup_{z\in[x,y]}\|\nabla f(z)\|)\|x-y\|$ globally; $\psi$ is measurable and convex on $Q$, with an exact measurable proximal oracle. IID samples from $\mu$ give a jointly measurable, integrable oracle with $\int g(x,\omega)\,d\mu=\nabla f(x)$ and $\int\|g(x,\omega)-\nabla f(x)\|^\alpha d\mu\le\sigma^\alpha$, $1<\alpha\le2$. A minimizer $x_\star$ satisfies $\|x_0-x_\star\|\le R$, $F(x_0)-F(x_\star)\le\Delta$, and $\|\nabla f(x_\star)\|\le G$, with $x_0\in Q$. Here $L_0,R,\Delta,\varepsilon>0$, $L_1,G,\sigma\ge0$, $0<\delta<1$. Target: $P\{x_{\rm out}\notin Q\ \text{or}\ F(x_{\rm out})-F(x_\star)>\varepsilon\}\le\delta$.}\par\smallskip
\begin{threeparttable}
\footnotesize
\setlength{\tabcolsep}{4pt}\renewcommand{\arraystretch}{1.25}
\begin{tabularx}{\linewidth}{@{}>{\raggedright\arraybackslash}p{0.20\linewidth}>{\raggedright\arraybackslash}p{0.32\linewidth}>{\raggedright\arraybackslash}p{0.25\linewidth}>{\raggedright\arraybackslash}X@{}}
\toprule
\textbf{Result} & \textbf{Oracle complexity} & \textbf{Novelty} & \textbf{Limitations} \\
\midrule
\textbf{(1)} \href{https://github.com/k3nfalt/longtails/blob/main/LongTails/Theorems/StochasticCompositeGeneralizedAccelerated.lean}{Stochastic\allowbreak Composite\allowbreak Generalized\allowbreak Accelerated} \textbf{(Ours)} $(\star\star\star)$ \href{https://github.com/k3nfalt/longtails/blob/main/scripts/reviews/StochasticCompositeGeneralizedAccelerated.md}{(review)} & $\mathcal{O}\!\left(2^{128\alpha/(\alpha-1)}\left[1+\sqrt{(L_0+L_1G)R^2/\varepsilon}+T+(\sigma R/\varepsilon)^p\right]H^{p+1}\right)$; gradient samples plus prox calls. & & \\
\midrule[1.2pt]
\textbf{(2)} \citep{scgsCompositeV2}, Thm.~2.7, $n=1$; derived ball reduction in $\mathbb R^d$ & $\mathcal{O}\!\left(1+\sqrt{L_bR^2/\varepsilon}\,h_b^2+\sqrt{RG/\varepsilon}\,h_b^{3/2}+(\sigma R/\varepsilon)^p h_b\right)$; samples plus prox calls. & \textbf{(1)}: polynomial rather than exponential curvature cost when that term dominates noise and logarithmic overhead. At $L_1=0$, removes $\sqrt{RG/\varepsilon}$. & \textbf{(1)} needs global smoothness and moments, rather than ball-local conditions; higher logarithmic powers and constants diverging as $\alpha\downarrow1$. \\
\midrule
\textbf{(3)} \citep{scgsOnePhaseV2}, Thm.~2, Remark~1; derived bound & $\mathcal{O}\!\left(h\left[1+\sqrt{L_0R^2/\varepsilon}+(\sigma R/\varepsilon)^p+(1+L_1R)\log\!\left(2+h(1+L_1R)\sqrt{\Delta/(L_0R^2)}\right)\right]\right)$; gradient samples. & \textbf{(1)} extends unconstrained acceleration to composite/constrained objectives with possibly nonzero $\nabla f(x_\star)$. \textbf{(3)} requires $Q=\mathbb R^d$, $\psi=0$, $G=0$, $L_1>0$. & \textbf{(1)} loses roughly $\log^p(\cdot/\delta)$ and an $\alpha$-dependent constant on the overlap; no rate or adaptivity gain. \\
\midrule
\textbf{(4)} \citep{scgsRepositoryFV} & $\mathcal{O}\!\left(\left[1+\sqrt{L_0R^2/\varepsilon}+D+(L_1RD)^{2/3}+\frac{\sigma^2R^2}{\varepsilon^2}(1+\log(1/\delta))^2\right]h_v\right)$; gradient samples. & \textbf{(1)} adds composite/constrained objectives and $1<\alpha<2$. \textbf{(4)} applies only to the unconstrained finite-variance case, $\psi=0$, $G=0$, $\alpha=2$. & \textbf{(1)} has a linear $L_1R$ transient instead of power $2/3$, and larger logarithmic overhead; no improvement on this overlap. \\
\bottomrule
\end{tabularx}
\begin{tablenotes}[flushleft]\footnotesize
\item[] \textbf{Abbr:} $p=\alpha/(\alpha-1)$; $T=(1+L_1R)\log(2+(1+L_1R)\sqrt{\Delta/(L_0R^2)})$; $H=\log((8+L_1R+(L_0+L_1G+L_1\sigma)R^2/\varepsilon+\Delta/(L_0R^2)+\sigma R/\varepsilon)/\delta)$; $L_b=(L_0+L_1G)e^{5\sqrt2L_1R}$; $h_b=1+\log((2+\sqrt{L_bR^2/\varepsilon}+\sqrt{RG/\varepsilon}+(\sigma R/\varepsilon)^p)/\delta)$; $h=1+\log(K/\delta)$ for the chosen horizon $K$ in \textbf{(3)}; $D=\log(2+\Delta/\varepsilon)$; $h_v=\log((2+L_0R^2/\varepsilon+L_1R+D)/\delta)$. $G$ enters only the analysis of \textbf{(1)}.
\end{tablenotes}
\end{threeparttable}
\end{table}
\clearpage
